\documentclass[11pt,a4paper]{article}
\usepackage[utf8]{inputenc}
\usepackage[vietnamese]{babel}
\usepackage{times}
\usepackage{geometry}
\geometry{left=3cm,right=2cm,top=2cm,bottom=2cm}
\usepackage{hyperref}
\usepackage{longtable}
\usepackage{booktabs}
\usepackage{amsmath}
\usepackage{graphicx}
\usepackage{indentfirst}
\usepackage[normalem]{ulem}
\usepackage{float}
\usepackage{color}
\usepackage{calc}
\usepackage{array}

% Định nghĩa tightlist để tránh lỗi biên dịch Pandoc
\providecommand{\tightlist}{%
  \setlength{\itemsep}{0pt}\setlength{\parskip}{0pt}}

% Cấu hình bảng biểu thụt lề và giãn dòng
\renewcommand{\arraystretch}{1.2}

% Tắt đánh số tự động cho các đề mục (do đã có đánh số thủ công trong nội dung)
\setcounter{secnumdepth}{-2}

\begin{document}

\noindent
\begin{minipage}[t]{0.45\textwidth}
\begin{center}
    \fontsize{11pt}{13pt}\selectfont TRƯỜNG ĐẠI HỌC BÁCH KHOA \\
    \vspace{0.15cm}
    \rule{3.2cm}{0.5pt}
\end{center}
\end{minipage}
\hfill
\begin{minipage}[t]{0.5\textwidth}
\begin{center}
    \fontsize{11pt}{13pt}\selectfont \textbf{CỘNG HÒA XÃ HỘI CHỦ NGHĨA VIỆT NAM} \\
    \vspace{0.1cm}
    \fontsize{11pt}{13pt}\selectfont \textbf{Độc lập - Tự do - Hạnh phúc} \\
    \vspace{0.05cm}
    \rule{4cm}{0.5pt}
\end{center}
\end{minipage}

\vspace{0.4cm}
\begin{flushright}
    \textit{Thành phố Hồ Chí Minh, ngày 11 tháng 6 năm 2026}
\end{flushright}

\vspace{0.6cm}
\begin{center}
    \fontsize{12pt}{14pt}\selectfont \textbf{KẾ HOẠCH TỔ CHỨC NGHIÊN CỨU KHOA HỌC} \\
    \vspace{0.15cm}
    \fontsize{12pt}{14pt}\selectfont \textbf{VÀ ĐỔI MỚI SÁNG TẠO CẤP SINH VIÊN 2026}
\end{center}

\vspace{0.5cm}

\subsection{I. TỔNG QUAN VÀ KẾ HOẠCH TRIỂN KHAI CỦA CÁC
KHOA}\label{i.-tux1ed5ng-quan-vuxe0-kux1ebf-houx1ea1ch-triux1ec3n-khai-cux1ee7a-cuxe1c-khoa}

Dưới đây là tóm tắt định hướng chi tiết và lộ trình triển khai hoạt động
Nghiên cứu Khoa học và Đổi mới Sáng tạo (NCKH \& ĐMST) cấp sinh viên của
10 khoa chuyên môn thuộc trường trong chu kỳ 12 tháng (từ tháng 6/2026
đến tháng 5/2027):

\textbf{1. Khoa Cơ khí}

\begin{itemize}
\tightlist
\item
  \textbf{Mục tiêu và Tờ trình}: Thực hiện theo Tờ trình số 151/TTr-ĐHBK
  ngày 05/06/2026, Khoa triển khai kế hoạch thí điểm 12 tháng nhằm thiết
  lập quy trình tổ chức NCKH sinh viên chuẩn hóa, làm tiền đề nhân rộng
  và lặp lại hiệu quả từ năm 2027.
\item
  \textbf{Định hướng chuyên môn}: Khoa tập trung huy động đông đảo lực
  lượng giảng viên hướng dẫn để hình thành các nhóm nghiên cứu nòng cốt,
  tạo tính kế thừa sản phẩm giữa các thế hệ sinh viên thay vì triển khai
  rời rạc. Kết quả nghiên cứu được hệ thống hóa làm minh chứng phục vụ
  kiểm định chất lượng đào tạo và làm nguồn sản phẩm tham gia các cuộc
  thi học thuật lớn.
\item
  \textbf{Lộ trình triển khai}: Tiến trình được chia làm 02 đợt. Đợt 1
  tiếp nhận và xét duyệt đề tài trong tháng 7/2026 (nghiệm thu vào tháng
  10/2026); Đợt 2 xét duyệt vào tháng 11-12/2026 (nghiệm thu tháng
  4/2027) và thực hiện thanh lý quyết toán toàn bộ kinh phí vào tháng
  5/2027.
\end{itemize}

\textbf{2. Khoa Kỹ thuật Xây dựng}

\begin{itemize}
\tightlist
\item
  \textbf{Mục tiêu và Tờ trình}: Thực hiện theo Tờ trình số 178/KTXD
  ngày 15/06/2026, Khoa hướng tới thúc đẩy mạnh mẽ phong trào NCKH và
  ĐMST trên tất cả các chương trình đào tạo của Khoa (Tiêu chuẩn, Kỹ sư
  tài năng - KSTN, PFIEV, và Đào tạo Quốc tế OISP).
\item
  \textbf{Lộ trình triển khai}: Quy trình thực hiện bám sát kế hoạch của
  Phòng KH\&CN: Thông báo đăng ký và tiếp nhận đề xuất của sinh viên
  trong tháng 6-7/2026; Thành lập Hội đồng xét duyệt cấp Khoa trong
  tháng 7/2026; Tổng hợp danh mục vật tư gửi Phòng KH\&CN mua sắm trong
  tháng 7-8/2026; Ký hợp đồng triển khai từ tháng 8-9/2026. Nghiệm thu
  đề tài vào tháng 4/2027 và quyết toán dứt điểm vào tháng 5/2027.
\end{itemize}

\textbf{3. Khoa Khoa học Ứng dụng}

\begin{itemize}
\tightlist
\item
  \textbf{Mục tiêu và Định hướng}: Khuyến khích tinh thần sáng tạo học
  thuật, phát triển năng lực nghiên cứu độc lập và khả năng ứng dụng
  thực tiễn của sinh viên. Khoa chủ trương chia các nhóm nghiên cứu sinh
  viên thành 2 hình thức: Nhóm A (Thực nghiệm chế tạo thiết bị, mô hình
  sinh học) và Nhóm B (Thuần túy mô phỏng tính toán lý thuyết).
\item
  \textbf{Hoạt động và Lộ trình}: Triển khai trong vòng 12 tháng theo
  quy trình chung của Nhà trường. Đặc biệt, Khoa dành riêng một phần
  kinh phí để tổ chức Hội thảo khoa học sinh viên cấp Khoa nhằm trưng
  bày sản phẩm thực tế, báo cáo kết quả học thuật và tạo diễn đàn trao
  đổi chuyên môn giữa sinh viên, giảng viên và các phòng thí nghiệm.
\end{itemize}

\textbf{4. Khoa Kỹ thuật Giao thông}

\begin{itemize}
\tightlist
\item
  \textbf{Mục tiêu và Kế hoạch}: Thực hiện theo Kế hoạch số 49/KH-KTGT
  ngày 12/06/2026, Khoa ưu tiên các đề tài hướng đến sản phẩm đầu ra
  kiểm chứng trực quan bao gồm mô hình vật lý, drone, xe tự hành, thiết
  bị thử nghiệm thực tế hoặc các phần mềm ứng dụng chuyên ngành điều
  khiển giao thông thông minh.
\item
  \textbf{Lộ trình và Đối tượng}: Đối tượng tham gia là sinh viên tất cả
  các hệ, trong đó ưu tiên sinh viên đang làm học phần tốt nghiệp. Nhận
  hồ sơ đăng ký trước ngày 31/07/2026, xét duyệt trong tháng 8/2026,
  thực hiện nghiên cứu từ tháng 11/2026 đến tháng 11/2027, nghiệm thu
  cấp Khoa vào tháng 10/2027 và thanh lý hợp đồng vào ngày 15/11/2027.
\end{itemize}

\textbf{5. Khoa Kỹ thuật Địa chất \& Dầu khí}

\begin{itemize}
\tightlist
\item
  \textbf{Mục tiêu và Định hướng}: Tập trung phát triển các đề tài
  nghiên cứu thực nghiệm chuyên sâu về cơ học đá, cơ học đất, địa kỹ
  thuật công trình và khảo sát thực địa đặc thù ngành. Cam kết đầu ra
  chất lượng cao là các công bố quốc tế uy tín (tạp chí thuộc danh mục
  Scopus Q3).
\item
  \textbf{Lộ trình triển khai}: Triển khai theo chu kỳ 6-12 tháng. Để
  đảm bảo tiến độ nghiên cứu thực nghiệm tại phòng thí nghiệm, Khoa ưu
  tiên tổng hợp và thực hiện quy trình mua sắm các thiết bị đo lường
  chuyên dụng ngay từ đợt đầu tiên.
\end{itemize}

\textbf{6. Khoa Công nghệ Vật liệu}

\begin{itemize}
\tightlist
\item
  \textbf{Mục tiêu và Định hướng}: Khuyến khích các đề tài nghiên cứu
  chế tạo vật liệu mới ứng dụng công nghệ cao và thân thiện với môi
  trường. Các hướng nghiên cứu chủ đạo gồm chế tạo vật liệu nano oxit
  kim loại làm anode cho pin lithium-ion, bê tông tự liền trong môi
  trường đặc thù, vật liệu in 3D từ nhựa tái chế và chế tạo khuôn
  concave vi mô.
\item
  \textbf{Lộ trình triển khai}: Triển khai chu kỳ 12 tháng. Khoa thực
  hiện rà soát và lập kế hoạch mua sắm hóa chất, vật tư thí nghiệm
  nghiêm ngặt từ sớm để tránh tình trạng chậm trễ mẫu thực nghiệm của
  sinh viên.
\end{itemize}

\textbf{7. Khoa Môi trường và Tài nguyên}

\begin{itemize}
\tightlist
\item
  \textbf{Mục tiêu và Định hướng}: Tập trung vào các nghiên cứu ứng dụng
  công nghệ hóa-sinh hiện đại giải quyết các thách thức môi trường thực
  tế. Các đề tài tiêu biểu gồm ứng dụng Fenton điện hóa và hoạt hóa
  peroxymonosulfate xử lý nước thải dệt nhuộm, khảo sát khả năng chịu
  mặn và tương thích Bacillus velezensis và Trichoderma asperellum hướng
  đến phát triển chế phẩm sinh học phục vụ nông nghiệp sạch.
\item
  \textbf{Lộ trình triển khai}: Triển khai từ tháng 08/2026 đến tháng
  07/2027. Khoa nỗ lực tối ưu hóa phân bổ nguồn kinh phí thực tế của nhà
  trường để đáp ứng nhu cầu thực hiện đề tài rất lớn của sinh viên.
\end{itemize}

\textbf{8. Khoa Khoa học \& Kỹ thuật Máy tính}

\begin{itemize}
\tightlist
\item
  \textbf{Mục tiêu và Định hướng}: Thúc đẩy nghiên cứu khoa học sinh
  viên đỉnh cao gắn liền với các công bố khoa học uy tín quốc tế và tham
  gia giải thưởng học thuật lớn. Định hướng đầu ra là bài báo đăng tại
  hội nghị có phản biện thuộc IEEE/Springer/ACM, Scopus Index hoặc tạp
  chí tính điểm của Hội đồng Giáo sư Nhà nước.
\item
  \textbf{Lộ trình triển khai}: Xét duyệt đề tài và mua sắm thiết bị
  phần cứng/phần mềm trong tháng 8/2026, thực hiện nghiên cứu từ tháng
  9/2026 đến tháng 4/2027, nghiệm thu và thanh lý vào tháng 5/2027.
  Khuyến khích sinh viên tham gia Huawei ICT Competition, Robocon, BK
  Innovation, Liên hoan Tuổi trẻ Sáng tạo.
\end{itemize}

\textbf{9. Khoa Quản lý Công nghiệp}

\begin{itemize}
\tightlist
\item
  \textbf{Mục tiêu và Định hướng}: Tập trung vào các đề tài ứng dụng
  công nghệ số và phương pháp tối ưu hóa trong quản lý kinh tế và vận
  hành chuỗi cung ứng. Hướng nghiên cứu chủ đạo gồm ứng dụng học máy
  trong dự báo nhu cầu doanh nghiệp, tối ưu hóa chuỗi cung ứng logistics
  (Heuristics, NSGA-II), định lượng khả năng phục hồi chuỗi cung ứng
  dưới tác động ripple effect, và ứng dụng AI trong lập kế hoạch du lịch
  xanh bền vững (mô hình UTAUT2).
\item
  \textbf{Lộ trình triển khai}: Triển khai chu kỳ 9-12 tháng bắt đầu từ
  tháng 9/2026, nghiệm thu và đánh giá kết quả theo đúng quy trình chung
  của Khoa và Nhà trường.
\end{itemize}

\textbf{10. Khoa Kỹ thuật Hóa học}

\begin{itemize}
\tightlist
\item
  \textbf{Mục tiêu và Tờ trình}: Thực hiện theo Tờ trình số 182/TTr-KTHH
  ngày 11/06/2026, Khoa thực hiện kế hoạch 12 tháng (tháng 6/2026 -
  tháng 5/2027) nhằm xây dựng quy trình quản lý NCKH sinh viên chuyên
  nghiệp và hiệu quả. Nâng cao số lượng sinh viên tiếp cận nghiên cứu
  thực nghiệm hóa chất, thúc đẩy công bố khoa học.
\item
  \textbf{Lộ trình}: Chia làm 2 đợt tiếp nhận và nghiệm thu tương tự
  Khoa Cơ khí, đảm bảo giải ngân và quyết toán dứt điểm toàn bộ kinh phí
  vào tháng 5/2027.
\end{itemize}

\begin{center}\rule{0.5\linewidth}{0.5pt}\end{center}

\subsection{II. THỐNG KÊ SỐ LƯỢNG ĐỀ TÀI VÀ SINH VIÊN THEO CHƯƠNG
TRÌNH}\label{ii.-thux1ed1ng-kuxea-sux1ed1-lux1b0ux1ee3ng-ux111ux1ec1-tuxe0i-vuxe0-sinh-viuxean-theo-chux1b0ux1a1ng-truxecnh}

Bảng dưới đây thống kê chi tiết số lượng đề tài, số lượng sinh viên tham
gia phân bổ theo 4 chương trình đào tạo chính (Tiêu chuẩn, Kỹ sư tài
năng - KSTN, Kỹ sư chất lượng cao Việt-Pháp - PFIEV, Đào tạo Quốc tế -
OISP) và tổng kinh phí của các khoa:

\begin{center}
\scriptsize
  \setlength{\arrayrulewidth}{0.4pt}
\begin{tabular}{|
  >{\centering\arraybackslash}p{(\linewidth - 22\tabcolsep) * \real{0.1136}}|
  >{\raggedright\arraybackslash}p{(\linewidth - 22\tabcolsep) * \real{0.0909}}|
  >{\centering\arraybackslash}p{(\linewidth - 22\tabcolsep) * \real{0.1136}}|
  >{\centering\arraybackslash}p{(\linewidth - 22\tabcolsep) * \real{0.1136}}|
  >{\centering\arraybackslash}p{(\linewidth - 22\tabcolsep) * \real{0.1136}}|
  >{\centering\arraybackslash}p{(\linewidth - 22\tabcolsep) * \real{0.1136}}|
  >{\centering\arraybackslash}p{(\linewidth - 22\tabcolsep) * \real{0.1136}}|
  >{\centering\arraybackslash}p{(\linewidth - 22\tabcolsep) * \real{0.1136}}|
  >{\centering\arraybackslash}p{(\linewidth - 16\tabcolsep) * \real{0.1136}}|}
\hline
\begin{minipage}[b]{\linewidth}\centering
STT
\end{minipage} & \begin{minipage}[b]{\linewidth}\raggedright
Tên Khoa
\end{minipage} & \begin{minipage}[b]{\linewidth}\centering
Chương trình Tiêu chuẩn (Đề tài / SV)
\end{minipage} & \begin{minipage}[b]{\linewidth}\centering
Chương trình KSTN (Đề tài / SV)
\end{minipage} & \begin{minipage}[b]{\linewidth}\centering
Chương trình PFIEV (Đề tài / SV)
\end{minipage} & \begin{minipage}[b]{\linewidth}\centering
Chương trình OISP (Đề tài / SV)
\end{minipage} & \begin{minipage}[b]{\linewidth}\centering
Tổng số Đề tài
\end{minipage} & \begin{minipage}[b]{\linewidth}\centering
Tổng số Sinh viên
\end{minipage} & \begin{minipage}[b]{\linewidth}\centering
Tổng Kinh phí (Triệu VNĐ)
\end{minipage} \\ \hline
\hline
1 & Khoa Cơ khí & 45-55 / 160-250 & 0 / 0 & 0 / 0 & 0 / 0 & \textbf{45 -
55} & \textbf{160 - 250} & \textbf{1.118,80} \\ \hline
2 & Kỹ thuật Xây dựng & 20 / 60 & 5 / 15 & 10 / 30 & 10 / 30 &
\textbf{45} & \textbf{135} & \textbf{440,60} \\ \hline
3 & Khoa học Ứng dụng & 4 / 16 & 13 / 18 & 0 / 0 & 4 / 9 & \textbf{21} &
\textbf{43} & \textbf{168,40} \\ \hline
4 & Kỹ thuật Giao thông & 10 / - & 0 / 0 & 10 / - & 10 / - & \textbf{30}
& \textbf{140} & \textbf{152,60} \\ \hline
5 & Kỹ thuật Địa chất \& Dầu khí & 3 / 22 & 0 / 0 & 0 / 0 & 0 / 0 &
\textbf{3} & \textbf{22} & \textbf{263,55} \\ \hline
6 & Công nghệ Vật liệu & 11 / 11 & 0 / 0 & 6 / 12 & 2 / 5 & \textbf{19}
& \textbf{28} & \textbf{94,70} \\ \hline
7 & Môi trường và Tài nguyên & 6 / 14 & 0 / 0 & 0 / 0 & 4 / 10 &
\textbf{10} & \textbf{24} & \textbf{27,40} \\ \hline
8 & Khoa học \& KT Máy tính & - / 70-100 & - / - & - / - & - / - &
\textbf{-} & \textbf{70 - 100} & \textbf{1.096,50} \\ \hline
9 & Quản lý Công nghiệp & 11 / 31 & 7 / 16 & 0 / 0 & 14 / 24 &
\textbf{32} & \textbf{71} & \textbf{Chưa rõ} \\ \hline
10 & Kỹ thuật Hóa học & 30-50 / 200-300 & 0 / 0 & 0 / 0 & 0 / 0 &
\textbf{30 - 50} & \textbf{200 - 300} & \textbf{848,70} \\ \hline
& \textbf{TỔNG CỘNG} & \textbf{140+ / 574+} & \textbf{25+ / 49} &
\textbf{26 / 72} & \textbf{44 / 78} & \textbf{235+} & \textbf{973+} &
\textbf{4.311,25} \\ \hline
\end{tabular}
\end{center}

\emph{Ghi chú}: * Các ô có ký tự \texttt{-} biểu thị thông tin chưa được
phân rã chi tiết trong kế hoạch gửi về của khoa. * Tổng kinh phí chưa
tính kinh phí của Khoa Quản lý Công nghiệp (do khoa chưa gửi con số phân
bổ cụ thể). * Số lượng đề tài/sinh viên của Khoa Cơ khí và Kỹ thuật Hóa
học được tính theo định biên mục tiêu tối thiểu.

\begin{center}\rule{0.5\linewidth}{0.5pt}\end{center}

\subsection{III. KẾ HOẠCH MUA SẮM NGUYÊN VẬT LIỆU, LINH KIỆN VÀ HÓA
CHẤT}\label{iii.-kux1ebf-houx1ea1ch-mua-sux1eafm-nguyuxean-vux1eadt-liux1ec7u-linh-kiux1ec7n-vuxe0-huxf3a-chux1ea5t}

Các đề tài NCKH sinh viên phần lớn đều có nhu cầu mua sắm vật tư, linh
kiện, hóa chất phục vụ cho thực nghiệm và chế tạo mô hình. Dưới đây là
tổng hợp nhu cầu kinh phí và danh mục mua sắm tiêu biểu:

\subsubsection{1. Tổng hợp kinh phí mua sắm theo đề xuất của các
khoa}\label{tux1ed5ng-hux1ee3p-kinh-phuxed-mua-sux1eafm-theo-ux111ux1ec1-xuux1ea5t-cux1ee7a-cuxe1c-khoa}

\begin{itemize}
\tightlist
\item
  \textbf{Khoa Cơ khí}: Dự kiến dành \textbf{670 triệu đồng} (chiếm 60\%
  tổng kinh phí được phân bổ) cho công tác mua sắm vật tư, linh kiện và
  gia công chế tạo.
\item
  \textbf{Khoa Khoa học \& Kỹ thuật Máy tính}: Dự kiến dành \textbf{596
  triệu đồng} (chiếm 54\% tổng kinh phí) cho mua sắm thiết bị phần cứng,
  linh kiện điện tử và bản quyền phần mềm nghiên cứu.
\item
  \textbf{Khoa Kỹ thuật Hóa học}: Dự kiến dành \textbf{500 triệu đồng}
  (chiếm 60\% tổng kinh phí) cho hóa chất tinh khiết, dụng cụ thủy tinh
  phòng thí nghiệm và thiết bị phụ trợ.
\item
  \textbf{Khoa Kỹ thuật Địa chất \& Dầu khí}: Đăng ký danh mục mua sắm
  chi tiết trị giá \textbf{263,55 triệu đồng}.
\item
  \textbf{Khoa Công nghệ Vật liệu}: Đăng ký danh mục mua sắm chi tiết
  trị giá \textbf{112,05 triệu đồng} (trong tổng đăng ký nhu cầu thực tế
  là 112,05 triệu đồng).
\item
  \textbf{Khoa Môi trường và Tài nguyên}: Đăng ký danh mục mua sắm chi
  tiết trị giá \textbf{67,27 triệu đồng} (Trường duyệt phân bổ thực tế
  là 27,4 triệu đồng).
\end{itemize}

\subsubsection{2. Danh mục mua sắm chi tiết tiêu biểu của các nhóm
nghiên
cứu}\label{danh-mux1ee5c-mua-sux1eafm-chi-tiux1ebft-tiuxeau-biux1ec3u-cux1ee7a-cuxe1c-nhuxf3m-nghiuxean-cux1ee9u}

\begin{itemize}
\tightlist
\item
  \textbf{Vật tư cơ khí, điện tử và tự động hóa (Khoa Cơ khí, Giao
  thông, Máy tính)}:

  \begin{itemize}
  \tightlist
  \item
    Linh kiện điện tử chuyên dụng: Vi điều khiển nhúng trên UAV, ESP32,
    mạch Arduino, cảm biến đo lực, cảm biến khoảng cách, cảm biến IMU,
    module truyền thông, module GPS, camera nhận diện hình ảnh.
  \item
    Vật liệu cơ khí chế tạo: Nhôm định hình, thép tấm, nhựa in 3D chuyên
    dụng, sợi carbon, vật liệu composite chịu lực cao dùng chế tạo vỏ
    khung xe tự hành, cánh sóng và các cụm thử nghiệm robot.
  \end{itemize}
\item
  \textbf{Hóa chất và vật liệu sinh học (Khoa Kỹ thuật Hóa học, Công
  nghệ Vật liệu, Môi trường, Khoa học Ứng dụng)}:

  \begin{itemize}
  \tightlist
  \item
    Hóa chất phân tích \& tổng hợp xúc tác: Hóa chất tinh khiết KOH,
    PVA, NaOH, H$_2$SO$_4$, HNO$_3$, cồn công nghiệp 96-98\%, dung dịch Nafion
    1100EW, các muối mangan(II) clorua MnCl2, thiếc(IV) clorua, đồng
    sunfat, sắt sunfat, boric acid, natri thiosunfat, amoni molybdat cấp
    AR.
  \item
    Thiết bị phụ trợ và dụng cụ thủy tinh: Điện cực glassy carbon chữ L
    6mm, micropipette kích thước 0.5-10 uL và 100-1000uL, hộp đựng tip
    nhựa polypropylene kháng hóa chất, chén nung sứ và nung kim loại
    100cc, màng bọc cao su chịu lực, giấy lọc, dây o-ring làm kín.
  \end{itemize}
\item
  \textbf{Thành phần thiết bị đo lường và khảo sát (Khoa Địa chất \& Dầu
  khí, Khoa học Ứng dụng)}:

  \begin{itemize}
  \tightlist
  \item
    Thiết bị quang học đo biến dạng: Ống kính macro chuyên dụng lắp trên
    camera đo biến dạng mẫu đá, đèn flash công suất lớn, tripod chịu tải
    cao.
  \item
    Thiết bị y sinh và y học: Gel tẩy da NuPrep, gel kết nối điện cực đo
    điện não Ten20 phục vụ phân tích tín hiệu EEG, thiết bị đo sinh tồn
    (PPG).
  \end{itemize}
\end{itemize}

\begin{center}\rule{0.5\linewidth}{0.5pt}\end{center}

\subsection{IV. ĐỊNH HƯỚNG CÔNG BỐ KHOA HỌC VÀ BÁO CÁO HỘI NGHỊ, HỘI
THẢO}\label{iv.-ux111ux1ecbnh-hux1b0ux1edbng-cuxf4ng-bux1ed1-khoa-hux1ecdc-vuxe0-buxe1o-cuxe1o-hux1ed9i-nghux1ecb-hux1ed9i-thux1ea3o}

Công bố khoa học và tham gia các cuộc thi học thuật là sản phẩm đầu ra
bắt buộc của hoạt động NCKH \& ĐMST sinh viên năm 2026:

\subsubsection{1. Định hướng xuất bản bài báo khoa
học}\label{ux111ux1ecbnh-hux1b0ux1edbng-xuux1ea5t-bux1ea3n-buxe0i-buxe1o-khoa-hux1ecdc}

\begin{itemize}
\tightlist
\item
  \textbf{Khoa KH\&KT Máy tính}: Định hướng bài báo được chấp nhận đăng
  tại các hội thảo/hội nghị uy tín thuộc danh mục IEEE, Springer, ACM;
  tạp chí thuộc danh mục Scopus hoặc tạp chí chuyên ngành được Hội đồng
  Chức danh Giáo sư Nhà nước (HDCDGSNN) tính điểm.
\item
  \textbf{Khoa Kỹ thuật Địa chất \& Dầu khí}: Cam kết đầu ra gồm ít nhất
  01 bài báo tạp chí quốc tế thuộc danh mục Scopus Q3 và 01 báo cáo tại
  Hội nghị khoa học quốc tế.
\item
  \textbf{Khoa Quản lý Công nghiệp}: Định hướng bài báo đăng tại các tạp
  chí trong danh mục HDCDGSNN, kỷ yếu hội nghị quốc tế và tạp chí Scopus
  Q4.
\item
  \textbf{Khoa Công nghệ Vật liệu \& Kỹ thuật Hóa học}: Bài báo gửi đăng
  tạp chí chuyên ngành trong nước/quốc tế và báo cáo poster tại Ngày hội
  kỹ thuật cấp khoa.
\end{itemize}

\subsubsection{2. Định hướng tham gia các cuộc thi học thuật
lớn}\label{ux111ux1ecbnh-hux1b0ux1edbng-tham-gia-cuxe1c-cuux1ed9c-thi-hux1ecdc-thuux1eadt-lux1edbn}

\begin{itemize}
\tightlist
\item
  Định hướng đưa các sản phẩm nghiên cứu thực tế tham gia tranh tài tại
  các cuộc thi công nghệ và học thuật uy tín như Huawei ICT Competition,
  Robocon, BK Innovation, Giải thưởng Khoa học và Công nghệ dành cho
  sinh viên trong các cơ sở giáo dục Đại học của Bộ GD\&ĐT, và Liên hoan
  Tuổi trẻ Sáng tạo.
\end{itemize}

\subsubsection{3. Tổng hợp kinh phí dành cho Hội nghị, Hội thảo (HN,
HT)}\label{tux1ed5ng-hux1ee3p-kinh-phuxed-duxe0nh-cho-hux1ed9i-nghux1ecb-hux1ed9i-thux1ea3o-hn-ht}

\begin{itemize}
\tightlist
\item
  \textbf{Khoa KH\&KT Máy tính}: Dành riêng \textbf{410 triệu đồng}
  (\textasciitilde37\% tổng kinh phí) để tài trợ chi phí đăng bài báo,
  lệ phí hội nghị và chi phí đi lại, ăn ở cho sinh viên báo cáo.
\item
  \textbf{Khoa Kỹ thuật Hóa học}: Dành \textbf{265 triệu đồng}
  (\textasciitilde30\% tổng kinh phí) hỗ trợ sinh viên công bố khoa học
  và đăng ký tham dự hội nghị.
\item
  \textbf{Khoa Cơ khí}: Dành \textbf{112 triệu đồng}
  (\textasciitilde10\% tổng kinh phí) hỗ trợ chi phí đăng bài và hội
  thảo khoa học liên quan trực tiếp đến đề tài.
\item
  \textbf{Khoa Khoa học Ứng dụng}: Dành \textbf{27,5 triệu đồng} từ
  nguồn phân bổ để tổ chức riêng Hội thảo khoa học sinh viên tại khoa
  phục vụ báo cáo và triển lãm.
\end{itemize}

\begin{center}\rule{0.5\linewidth}{0.5pt}\end{center}

\subsection{V. TỔNG HỢP CÁC ĐỀ XUẤT VÀ KIẾN NGHỊ CỦA CÁC
KHOA}\label{v.-tux1ed5ng-hux1ee3p-cuxe1c-ux111ux1ec1-xuux1ea5t-vuxe0-kiux1ebfn-nghux1ecb-cux1ee7a-cuxe1c-khoa}

Để hoạt động NCKH \& ĐMST cấp sinh viên năm 2026 đạt hiệu quả thực chất,
các khoa đã đưa ra một số đề xuất quan trọng gửi Phòng Khoa học \& Công
nghệ, Phòng Kế hoạch - Tài chính và Ban Giám hiệu Trường:

\subsubsection{1. Đề xuất về cơ chế tài chính và định mức hỗ trợ (Khoa
Cơ khí, Hóa học, Địa chất, Môi
trường)}\label{ux111ux1ec1-xuux1ea5t-vux1ec1-cux1a1-chux1ebf-tuxe0i-chuxednh-vuxe0-ux111ux1ecbnh-mux1ee9c-hux1ed7-trux1ee3-khoa-cux1a1-khuxed-huxf3a-hux1ecdc-ux111ux1ecba-chux1ea5t-muxf4i-trux1b0ux1eddng}

\begin{itemize}
\tightlist
\item
  \textbf{Hỗ trợ kinh phí sinh viên}: Kiến nghị áp dụng cơ chế cho phép
  chi bồi dưỡng trực tiếp cho sinh viên tham gia thực hiện đề tài (dự
  kiến định mức \(\approx\)15\% giá trị đề tài) nhằm tạo động lực và nâng cao
  trách nhiệm nghiên cứu của sinh viên.
\item
  \textbf{Đặc thù ngành}: Kiến nghị xây dựng định mức chi phí đi lại,
  khảo sát thực địa và thu thập số liệu hiện trường cho sinh viên (dự
  kiến chiếm \(\approx\)10\% kinh phí phân bổ), đặc biệt cần thiết cho các ngành
  có khối lượng công việc thực địa lớn như Địa chất Dầu khí, Môi trường
  và Tài nguyên.
\item
  \textbf{Tăng hạn mức kinh phí}: Một số khoa (như Môi trường và Tài
  nguyên, Công nghệ Vật liệu) đề nghị nhà trường xem xét nâng hạn mức
  kinh phí phân bổ thực tế vì nhu cầu đăng ký và tính toán dự toán hóa
  chất, thiết bị thực tế của sinh viên vượt xa hạn mức trường duyệt.
\end{itemize}

\subsubsection{2. Đề xuất cải tiến quy trình hành chính và mua sắm (Khoa
KHUD, Giao thông, Máy tính, Vật
liệu)}\label{ux111ux1ec1-xuux1ea5t-cux1ea3i-tiux1ebfn-quy-truxecnh-huxe0nh-chuxednh-vuxe0-mua-sux1eafm-khoa-khud-giao-thuxf4ng-muxe1y-tuxednh-vux1eadt-liux1ec7u}

\begin{itemize}
\tightlist
\item
  \textbf{Quy trình hóa chất \& thiết bị chuyên dụng}: Hướng dẫn chi
  tiết quy trình mua sắm, nhập khẩu và quản lý sử dụng các hóa chất thí
  nghiệm thuộc danh mục kiểm soát (hóa chất độc hại), các linh kiện điện
  tử chuyên dụng nhập khẩu và các phần mềm bản quyền mô phỏng nghiên
  cứu.
\item
  \textbf{Rút ngắn thời gian giải ngân}: Đề xuất tối ưu hóa quy trình
  duyệt hồ sơ mua sắm và giải ngân (rút ngắn thời gian thông thường 30
  ngày xuống mức tối thiểu) để tránh làm gián đoạn tiến độ thực hiện đề
  tài của sinh viên, đặc biệt với các đề tài chế tạo thực nghiệm chỉ có
  chu kỳ 12 tháng.
\end{itemize}

\subsubsection{3. Đề xuất cơ chế ghi nhận cho giảng viên hướng dẫn (Khoa
Cơ khí, KH\&KTMT,
KTHH)}\label{ux111ux1ec1-xuux1ea5t-cux1a1-chux1ebf-ghi-nhux1eadn-cho-giux1ea3ng-viuxean-hux1b0ux1edbng-dux1eabn-khoa-cux1a1-khuxed-khktmt-kthh}

\begin{itemize}
\tightlist
\item
  \textbf{Ghi nhận công sức giảng viên}: Đề nghị Phòng KH\&CN ban hành
  hướng dẫn cụ thể về việc tính điểm Nhiệm vụ 2 (quy đổi giờ chuẩn giảng
  dạy) và hỗ trợ kinh phí bồi dưỡng hướng dẫn cho giảng viên hướng dẫn
  đề tài NCKH sinh viên nhằm khuyến khích giảng viên đồng hành lâu dài
  cùng các nhóm sinh viên.
\end{itemize}

\subsubsection{4. Đề xuất cơ chế phối hợp tham gia cuộc thi ngoài trường
(Khoa Cơ khí,
KTHH)}\label{ux111ux1ec1-xuux1ea5t-cux1a1-chux1ebf-phux1ed1i-hux1ee3p-tham-gia-cuux1ed9c-thi-ngouxe0i-trux1b0ux1eddng-khoa-cux1a1-khuxed-kthh}

\begin{itemize}
\tightlist
\item
  \textbf{Tạo điều kiện tham gia cuộc thi}: Kiến nghị Phòng Công tác
  Sinh viên (CTSV) hỗ trợ đơn giản hóa quy trình tiếp nhận hồ sơ, làm
  thủ tục đề cử sinh viên tham gia các cuộc thi học thuật ngoài trường
  (theo quy trình BK-QT-ED-056-05) nhằm tạo cơ hội cọ xát thực tế cho
  sinh viên sau khi hoàn thành đề tài cấp cơ sở.
\end{itemize}


\end{document}
